% Capitulo 5: listados de codigo en los cuatro estilos, mas ink.
\section{Listados de código}\label{sec:codigo}

texforge resalta el código fuente durante la compilación: cada bloque
\texttt{code} se reescribe sobre la copia de compilación y solo necesita
\texttt{color.sty}, sin \texttt{listings} ni \texttt{minted}. El lenguaje se
indica con \texttt{lang} y la numeración con \texttt{numbers}. El tema de
color y el estilo por defecto se eligen una vez para todo el documento en
\texttt{project.toml}, aquí \texttt{one-light} con el estilo
\texttt{light}, y cada bloque puede pedir su propio estilo: hay
cuatro --- \texttt{light}, \texttt{light-mono}, \texttt{dark} y
\texttt{dark-mono} --- más el quinto estilo \texttt{ink}, en negro.

\subsection{Python}

El listado~\ref{lst:python} aproxima la integral de Gauss del
capítulo~\ref{sec:matematicas} por cuadratura: el mismo $\sqrt{\pi}$ de la
ecuación~\eqref{eq:gauss}, ahora calculado a fuerza bruta.

\begin{code}[lang=python, numbers=true, caption={Cuadratura de la integral de Gauss}, label={lst:python}]
import numpy as np

def integrar_gauss(n_puntos=1000):
    """Aproxima la integral de Gauss por cuadratura."""
    x = np.linspace(-10, 10, n_puntos)
    dx = x[1] - x[0]
    y = np.exp(-x**2)
    return np.sum(y) * dx

if __name__ == "__main__":
    print(f"Integral = {integrar_gauss():.6f}")
\end{code}

\subsection{Rust}

El listado~\ref{lst:fib} muestra la definición clásica de Fibonacci en
Rust con el estilo \texttt{light-mono}: sin color, las palabras clave van en
negrita y los comentarios en cursiva, que es lo que funciona impreso en
blanco y negro.

\begin{code}[lang=rust, style=light-mono, numbers=true, caption={Fibonacci en Rust}, label={lst:fib}]
fn fibonacci(n: u64) -> u64 {
    match n {
        0 => 0,
        1 => 1,
        _ => fibonacci(n - 1) + fibonacci(n - 2),
    }
}

fn main() {
    println!("fib(10) = {}", fibonacci(10));
}
\end{code}

\subsection{Shell}

El listado~\ref{lst:shell} usa el estilo \texttt{dark}, con el aire de una
terminal. Los caracteres especiales de \LaTeX{} se imprimen tal cual, porque
cada carácter se escapa antes de llegar al motor.

\begin{code}[lang=bash, style=dark, numbers=true, caption={Limpieza del directorio de compilación}, label={lst:shell}]
#!/usr/bin/env bash
set -euo pipefail

texforge build
texforge pdf info
texforge clean
\end{code}

\subsection{TOML}

El listado~\ref{lst:toml} es la configuración de resaltado de este
documento, con el estilo \texttt{dark-mono}: fondo oscuro y una sola tinta
clara, negritas y cursivas en lugar de colores.

\begin{code}[lang=toml, style=dark-mono, numbers=true, caption={Resaltado del propio documento}, label={lst:toml}]
[highlight]
theme = "one-light"
style = "light"
\end{code}

\subsection{Código desde un archivo}

El listado~\ref{lst:archivo} no está escrito en el \texttt{.tex}: se lee
del guion \texttt{scripts/stats.py} con \texttt{file} y un rango con
\texttt{lines}. Los números son los del archivo, y un archivo ausente o un
rango inválido fallan la compilación y \texttt{texforge check}.

\begin{code}[lang=python, file=scripts/stats.py, lines=5-24, numbers=true, caption={Funciones leídas del archivo}, label={lst:archivo}]
\end{code}

\subsection{Estilo ink}

El mismo archivo con \texttt{style=ink} usa solo negro: distingue los
elementos por peso y forma, con palabras clave y números en negrita,
nombres en negrita cursiva y comentarios en cursiva.

\begin{code}[lang=python, file=scripts/stats.py, lines=5-24, numbers=true, style=ink]
\end{code}

\begin{novabox}
  El estilo \texttt{ink} es nuevo en la 1.0: ignora el tema y compone en
  negro sobre blanco, pensado para imprimir o fotocopiar.
\end{novabox}

\subsection{Pies y etiquetas}

Todo bloque puede llevar pie y etiqueta: \texttt{caption=} escribe la línea
de Listado con su número y \texttt{label=} le da nombre para citarlo con
\texttt{\textbackslash ref}, como hacen los listados de este capítulo.
\texttt{pos=} decide si flota o se queda donde está escrito, y
\texttt{\textbackslash listoflistings} reúne todos los pies en un índice
propio, igual que el de figuras y tablas.

\subsection{Copiar sin números}

Los números de línea se omiten al copiar en los visores que respetan el
\texttt{ActualText} del PDF, y \texttt{numbers=false} los quita en todas
partes, como en el bloque siguiente.

\begin{code}[lang=python, file=scripts/stats.py, lines=5-24, numbers=false]
\end{code}

\begin{notebox}
  Pegar de un listado numerado no arrastra los números en los visores
  modernos: el PDF guarda el texto real aparte de lo que se ve. Si tu flujo
  exige páginas sin números, \texttt{numbers=false} los retira del todo.
\end{notebox}
